\documentclass[12pt]{scrartcl}
\usepackage[hagavi]{azzam}

% Fungsi Aritmetika - Modul (gabungan Materi dan Soal Latihan)
% Fungsi aritmetika: tau, sigma, phi, dan fungsi multiplikatif.
% 50 soal: 20 Banyak Faktor, 11 Jumlah Faktor, 4 Totient, 9 Campuran, 6 Tantangan.
% Adaptasi dari Bab 3 buku Modern Olympiad Number Theory (Aditya Khurmi)
% dan 8 Outsource K Math/Pra-OSP/FungsiAritmetika.tex.
% _Ujian dan _Solusi tetap terpisah (menyusul).
% Sumber tiap soal latihan dimunculkan di PDF sebagai \item (Sumber); komentar menyimpan detail lengkap.
% Soal pilihan ganda asli disajikan tanpa opsi jawaban sebagai soal isian.

\title{Fungsi Aritmetika - Modul}
\date{15 Oktober 2026}

\begin{document}
\maketitle

\renewcommand*\contentsname{Daftar Isi}
\tableofcontents

\newpage

\part*{Materi}

Modul ini membahas \istilah{fungsi aritmetika}{arithmetic function}, yaitu fungsi
yang mengaitkan setiap bilangan asli dengan sebuah nilai. Tiga yang paling sering
muncul di OSK dan OSP adalah $\tau$ (banyak faktor), $\sigma$ (jumlah faktor),
dan $\varphi$ (totient Euler, \textit{Euler's totient}).
Sifat kunci ketiganya sama: multiplikatif. Artinya nilai di perkalian dua bilangan
yang relatif prima (\textit{coprime}) sama dengan perkalian nilainya masing-masing.

\section{Fungsi Aritmetika dan Fungsi Multiplikatif}

\begin{definition}
Sebuah fungsi $f : \mathbb{N} \to \mathbb{C}$ disebut \textbf{fungsi aritmetika}.
Fungsi aritmetika $f$ disebut \istilah{multiplikatif}{multiplicative} apabila
\[ f(mn) = f(m)f(n) \]
untuk setiap $m$ dan $n$ yang relatif prima. Jika kesamaan itu berlaku untuk
\emph{semua} $m$ dan $n$, maka $f$ disebut \istilah{multiplikatif sempurna}{completely multiplicative}.
\end{definition}

\begin{fact}
\begin{enumerate}
    \item Fungsi $\tau$, $\sigma$, dan $\varphi$ semuanya multiplikatif
    (dibuktikan satu per satu di bagian masing-masing).
    \item Fungsi konstan $f(n) = 1$ dan fungsi identitas $\text{id}(n) = n$
    keduanya multiplikatif.
    \item Fungsi $g(n) = n^k$ multiplikatif sempurna untuk setiap bilangan cacah $k$.
    \item Fungsi multiplikatif ditentukan sepenuhnya oleh nilainya di perpangkatan
    prima (\textit{prime power}): untuk $n = p_1^{\alpha_1} \cdots p_k^{\alpha_k}$ berlaku
    $f(n) = f(p_1^{\alpha_1}) \cdots f(p_k^{\alpha_k})$.
\end{enumerate}
\end{fact}

\begin{example}
Diketahui $f$ multiplikatif dengan $f(8) = 5$ dan $f(9) = 7$.
Tentukan $f(72)$.
\end{example}
\begin{proof}[Solusi]
Karena $72 = 8 \cdot 9$ dengan $FPB(8, 9) = 1$:
\[ f(72) = f(8)f(9) = 5 \cdot 7 = 35. \]
\end{proof}

\begin{example}
Tunjukkan bahwa $g(n) = n^2$ multiplikatif sempurna.
\end{example}
\begin{proof}[Solusi]
Untuk setiap $m$ dan $n$:
\[ g(mn) = (mn)^2 = m^2n^2 = g(m)g(n). \]
\end{proof}

\section{Banyak Faktor Positif}

\begin{definition}
Untuk bilangan asli $n$, notasi $\tau(n)$ menyatakan banyak faktor positif
(\textit{number of divisors}) dari $n$.
\end{definition}

\begin{theorem}[Rumus $\tau$]
Jika $n = p_1^{\alpha_1} p_2^{\alpha_2} \cdots p_k^{\alpha_k}$ dengan $p_i$ prima berbeda, maka
\[ \tau(n) = (\alpha_1 + 1)(\alpha_2 + 1) \cdots (\alpha_k + 1). \]
\end{theorem}

Alasannya pencacahan langsung: setiap faktor $n$ berbentuk
$p_1^{\beta_1} \cdots p_k^{\beta_k}$ dengan $0 \le \beta_i \le \alpha_i$,
jadi ada $\alpha_i + 1$ pilihan untuk tiap $i$.

\begin{example}
Hitunglah $\tau(360)$.
\end{example}
\begin{proof}[Solusi]
Faktorkan dahulu:
\[ 360 = 2^3 \cdot 3^2 \cdot 5. \]
Maka
\[ \tau(360) = (3 + 1)(2 + 1)(1 + 1) = 4 \cdot 3 \cdot 2 = 24. \]
\end{proof}

\begin{fact}[$\tau$ ganjil]
$\tau(n)$ ganjil jika dan hanya jika $n$ \istilah{bilangan kuadrat}{perfect square}.
\end{fact}

Buktinya satu baris: $(\alpha_1 + 1) \cdots (\alpha_k + 1)$ ganjil tepat ketika
semua $\alpha_i$ genap, dan itu artinya $n$ kuadrat.

\begin{example}
Tentukan banyak bilangan asli $n \le 100$ dengan $\tau(n)$ ganjil.
\end{example}
\begin{proof}[Solusi]
$\tau(n)$ ganjil tepat ketika $n$ kuadrat.
Kuadrat yang tidak lebih dari $100$ adalah
\[ 1^2, 2^2, \dots, 10^2, \]
sebanyak $10$ bilangan.
\end{proof}

\begin{fact}[Hasil kali semua faktor]
Hasil kali semua faktor positif dari $n$ sama dengan $n^{\tau(n)/2}$.
\end{fact}

Pasangkan setiap faktor $d$ dengan pasangannya $n/d$.
Tiap pasangan berhasil kali $n$, dan ada $\tau(n)/2$ pasangan
(faktor tengah $\sqrt{n}$ berpasangan dengan dirinya sendiri bila $n$ kuadrat).

\begin{example}
Tentukan hasil kali semua faktor positif dari $12$.
\end{example}
\begin{proof}[Solusi]
Faktorisasi $12 = 2^2 \cdot 3$ memberi
\[ \tau(12) = (2 + 1)(1 + 1) = 6. \]
Hasil kalinya:
\[ 12^{6/2} = 12^3 = 1728. \]
\end{proof}

\begin{fact}[Batas $\tau$]
Untuk setiap bilangan asli $n$ berlaku $\tau(n) \le 2\sqrt{n}$.
\end{fact}

Tulis faktor dari $n$ secara berpasangan $(d, n/d)$.
Dari tiap pasangan, ambil anggota yang tidak lebih dari $\sqrt{n}$.
Pasangan berbeda memberi bilangan berbeda, dan banyak bilangan asli yang
tidak lebih dari $\sqrt{n}$ ada $\lfloor \sqrt{n} \rfloor \le \sqrt{n}$.
Jadi banyak pasangan tidak lebih dari $\sqrt{n}$, sehingga $\tau(n) \le 2\sqrt{n}$.

\begin{example}
Buktikan $\tau(n) \le 10$ untuk setiap $n \le 25$.
\end{example}
\begin{proof}[Solusi]
Batas di atas memberi
\[ \tau(n) \le 2\sqrt{n} \le 2\sqrt{25} = 10. \]
\end{proof}

\begin{lemma}[Menghitung ganda $\tau$]
Untuk setiap bilangan asli $N$ berlaku
\[ \tau(1) + \tau(2) + \dots + \tau(N) =
\left\lfloor \frac{N}{1} \right\rfloor + \left\lfloor \frac{N}{2} \right\rfloor
+ \dots + \left\lfloor \frac{N}{N} \right\rfloor. \]
\end{lemma}

Argumennya \istilah{menghitung ganda}{double counting}: kedua ruas menghitung
himpunan pasangan yang sama, yaitu $\{(d, k) : d \mid k,\ 1 \le k \le N\}$.
Ruas kiri menghitung per $k$ (ada $\tau(k)$ nilai $d$ untuk tiap $k$).
Ruas kanan menghitung per $d$: banyak kelipatan $d$ yang tidak lebih dari $N$
adalah $\lfloor N/d \rfloor$.

\begin{example}[Disederhanakan dari Tuymaada 2000]
Hitunglah $\tau(1) + \tau(2) + \dots + \tau(10)$ dengan dua cara.
\end{example}
\begin{proof}[Solusi]
Langsung dari definisi:
\[ 1 + 2 + 2 + 3 + 2 + 4 + 2 + 4 + 3 + 4 = 27. \]
Dengan lemma di atas:
\begin{align*}
\left\lfloor \frac{10}{1} \right\rfloor + \dots + \left\lfloor \frac{10}{10} \right\rfloor
&= 10 + 5 + 3 + 2 + 2 + 1 + 1 + 1 + 1 + 1 \\
&= 27.
\end{align*}
\end{proof}

\section{Jumlah Faktor Positif}

\begin{definition}
Untuk bilangan asli $n$, notasi $\sigma(n)$ menyatakan jumlah semua faktor
positif (\textit{sum of divisors}) dari $n$.
\end{definition}

\begin{theorem}[Rumus $\sigma$]
Fungsi $\sigma$ multiplikatif. Untuk $p$ prima dan $\alpha \ge 1$:
\[ \sigma(p^\alpha) = 1 + p + \dots + p^\alpha = \frac{p^{\alpha+1} - 1}{p - 1}. \]
Akibatnya, jika $n = p_1^{\alpha_1} \cdots p_k^{\alpha_k}$, maka
\[ \sigma(n) = \sigma(p_1^{\alpha_1}) \cdots \sigma(p_k^{\alpha_k}). \]
\end{theorem}

Alasannya pengelompokan: setiap faktor $n$ berbentuk $p_1^{\beta_1} \cdots p_k^{\beta_k}$,
dan jumlah atas semua pilihan $\beta_i$ memfaktorkan menjadi perkalian
jumlah per prima, persis seperti menguraikan $(a_1 + a_2)(b_1 + b_2)$.

\begin{example}
Hitunglah $\sigma(72)$.
\end{example}
\begin{proof}[Solusi]
Faktorisasi $72 = 2^3 \cdot 3^2$ memberi
\begin{align*}
\sigma(72) &= \sigma(2^3)\sigma(3^2) \\
&= \frac{2^4 - 1}{2 - 1} \cdot \frac{3^3 - 1}{3 - 1} \\
&= 15 \cdot 13 \\
&= 195.
\end{align*}
\end{proof}

\begin{definition}
Sebuah \istilah{bilangan sempurna}{perfect number} adalah bilangan asli $n$
dengan $\sigma(n) = 2n$.
\end{definition}

\begin{example}
Tunjukkan bahwa $28$ bilangan sempurna.
\end{example}
\begin{proof}[Solusi]
Faktorisasi $28 = 2^2 \cdot 7$ memberi
\begin{align*}
\sigma(28) &= \sigma(2^2)\sigma(7) \\
&= (1 + 2 + 4)(1 + 7) \\
&= 7 \cdot 8 \\
&= 56 = 2 \cdot 28.
\end{align*}
\end{proof}

\begin{theorem}[Euclid-Euler]
Bilangan asli genap $n$ sempurna jika dan hanya jika
$n = 2^{p-1}(2^p - 1)$ dengan $2^p - 1$ bilangan prima.
\end{theorem}

\begin{proof}[Solusi]
Untuk arah ($\Leftarrow$), jika $2^p - 1$ prima maka ia relatif prima dengan $2^{p-1}$,
sehingga
\begin{align*}
\sigma(2^{p-1}(2^p - 1)) &= \sigma(2^{p-1})\sigma(2^p - 1) \\
&= (2^p - 1) \cdot 2^p \\
&= 2n.
\end{align*}
Untuk arah ($\Rightarrow$), tulis $n = 2^k q$ dengan $q$ ganjil dan $k \ge 1$.
Karena $n$ sempurna:
\[ 2n = 2^{k+1}q = \sigma(n) = \sigma(2^k)\sigma(q) = (2^{k+1} - 1)\sigma(q). \]
Karena $2^{k+1} - 1$ ganjil, ia membagi $q$.
Tulis $q = x(2^{k+1} - 1)$, lalu substitusikan:
\[ 2^{k+1}x = \sigma(q). \]
Padahal $x \mid q$ dan $x < q$, jadi $\sigma(q) \ge q + x = 2^{k+1}x$.
Kesamaan memaksa $x = 1$: jika $x > 1$ maka $1$, $x$, $q$ tiga faktor berbeda
dan $\sigma(q) > q + x$. Jadi $q = 2^{k+1} - 1$ sekaligus $\sigma(q) = q + 1$,
artinya $q$ prima. Tulis $p = k + 1$ dan selesai.
\end{proof}

\begin{example}
Tentukan semua bilangan sempurna genap yang kurang dari $100$.
\end{example}
\begin{proof}[Solusi]
Periksa $2^{p-1}(2^p - 1) < 100$ dengan $2^p - 1$ prima:
\begin{align*}
p = 2 &: 2 \cdot 3 = 6, \\
p = 3 &: 4 \cdot 7 = 28, \\
p = 5 &: 16 \cdot 31 = 496 > 100.
\end{align*}
Untuk $p = 4$ nilai $2^4 - 1 = 15$ bukan prima.
Jadi jawabannya $6$ dan $28$.
\end{proof}

\begin{remark}
Bilangan prima berbentuk $2^p - 1$ disebut \istilah{prima Mersenne}{Mersenne prime}.
Sampai hari ini tidak diketahui apakah ada bilangan sempurna ganjil,
dan tidak diketahui apakah banyak bilangan sempurna tak hingga.
\end{remark}

\begin{fact}[$\sigma$ ganjil]
$\sigma(n)$ ganjil jika dan hanya jika $n$ bilangan kuadrat
atau dua kali bilangan kuadrat.
\end{fact}

Perhatikan $\sigma(p^\alpha) = 1 + p + \dots + p^\alpha$.
Jumlah ini ganjil tepat ketika $p$ ganjil dengan $\alpha$ genap,
atau $p = 2$ (dengan $\alpha$ berapa pun).
Jadi $\sigma(n)$ ganjil tepat ketika semua eksponen prima ganjil genap,
artinya $n = 2^a m^2$ dengan $m$ ganjil.
Jika $a$ genap maka $n$ kuadrat; jika $a$ ganjil maka $n$ dua kali kuadrat.

\begin{example}
Manakah yang ganjil, $\sigma(12)$ atau $\sigma(18)$?
\end{example}
\begin{proof}[Solusi]
Faktorisasi $12 = 2^2 \cdot 3^1$: eksponen $3$ ganjil, jadi $\sigma(12)$ genap
(cek: $1 + 2 + 3 + 4 + 6 + 12 = 28$).
Faktorisasi $18 = 2^1 \cdot 3^2$: dua kali kuadrat, jadi $\sigma(18)$ ganjil
(cek: $1 + 2 + 3 + 6 + 9 + 18 = 39$).
\end{proof}

\section{Fungsi Totient Euler}

\begin{definition}
Untuk bilangan asli $n$, notasi $\varphi(n)$ menyatakan banyak bilangan asli
$k \le n$ dengan $FPB(k, n) = 1$.
\end{definition}

\begin{example}
Tentukan $\varphi(12)$ dengan mendaftar satu per satu.
\end{example}
\begin{proof}[Solusi]
Bilangan dari $1$ sampai $12$ yang relatif prima dengan $12$ adalah
\[ 1,\ 5,\ 7,\ 11, \]
jadi $\varphi(12) = 4$.
\end{proof}

\begin{theorem}[$\varphi$ multiplikatif]
Jika $FPB(m, n) = 1$ maka $\varphi(mn) = \varphi(m)\varphi(n)$.
\end{theorem}

Susun $1, 2, \dots, mn$ dalam tabel $n$ baris dan $m$ kolom:
baris ke-$i$ berisi $(i-1)m + 1$ sampai $(i-1)m + m$.
Tiap baris adalah sistem residu lengkap (\textit{complete residue system}) modulo $m$,
jadi memuat $\varphi(m)$ bilangan yang relatif prima dengan $m$.
Tiap kolom adalah barisan aritmetika dengan beda $m$;
karena $FPB(m, n) = 1$, kolom itu sistem residu lengkap modulo $n$,
jadi memuat $\varphi(n)$ bilangan yang relatif prima dengan $n$.
Bilangan yang relatif prima dengan $mn$ adalah yang relatif prima
dengan $m$ sekaligus $n$: ada $\varphi(m)\varphi(n)$ buah.

\begin{example}
Diketahui $\varphi(7) = 6$ dan $\varphi(9) = 6$. Tentukan $\varphi(63)$.
\end{example}
\begin{proof}[Solusi]
Karena $63 = 7 \cdot 9$ dengan $FPB(7, 9) = 1$:
\[ \varphi(63) = \varphi(7)\varphi(9) = 6 \cdot 6 = 36. \]
\end{proof}

\begin{theorem}[Rumus $\varphi$]
Jika $n = p_1^{\alpha_1} \cdots p_k^{\alpha_k}$, maka
\[ \varphi(n) = n\left(1 - \frac{1}{p_1}\right) \cdots \left(1 - \frac{1}{p_k}\right). \]
\end{theorem}

Cukup periksa untuk perpangkatan prima lalu gunakan sifat multiplikatif:
$\varphi(p^\alpha) = p^\alpha - p^{\alpha-1}$ karena di antara $1, \dots, p^\alpha$
yang tidak relatif prima dengan $p^\alpha$ hanyalah $p^{\alpha-1}$ kelipatan $p$.

Ada juga bukti langsung memakai \istilah{prinsip inklusi eksklusi}{inclusion-exclusion principle},
tanpa memakai sifat multiplikatif. Pandang bilangan $1, 2, \dots, n$ dan untuk tiap $i$
misalkan $A_i$ himpunan kelipatan $p_i$ di antaranya. Bilangan yang dihitung $\varphi(n)$
tepat yang tidak masuk $A_i$ mana pun, jadi
\[ \varphi(n) = n - |A_1 \cup \dots \cup A_k|. \]
Inklusi eksklusi, yaitu $|A \cup B| = |A| + |B| - |A \cap B|$ yang diperluas
ke $k$ himpunan, memberi
\begin{align*}
|A_1 \cup \dots \cup A_k|
&= \sum_i |A_i| - \sum_{i<j} |A_i \cap A_j| + \dots \\
&= \sum_i \frac{n}{p_i} - \sum_{i<j} \frac{n}{p_ip_j} + \dots,
\end{align*}
karena $|A_{i_1} \cap \dots \cap A_{i_r}| = n/(p_{i_1} \cdots p_{i_r})$.
Maka
\[ \varphi(n) = n\left(1 - \sum_i \frac{1}{p_i} + \sum_{i<j} \frac{1}{p_ip_j}
- \dots\right) = n\left(1 - \frac{1}{p_1}\right) \cdots \left(1 - \frac{1}{p_k}\right). \]

\begin{example}
Hitunglah $\varphi(72)$.
\end{example}
\begin{proof}[Solusi]
Faktorisasi $72 = 2^3 \cdot 3^2$ memberi
\[ \varphi(72) = 72\left(1 - \frac{1}{2}\right)\left(1 - \frac{1}{3}\right)
= 72 \cdot \frac{1}{2} \cdot \frac{2}{3} = 24. \]
\end{proof}

\begin{theorem}[Gauss]
Untuk setiap bilangan asli $n$ berlaku
\[ \sum_{d \mid n} \varphi(d) = n. \]
\end{theorem}

Pandang $n$ pecahan $\frac{1}{n}, \frac{2}{n}, \dots, \frac{n}{n}$,
lalu sederhanakan masing-masing.
Penyebut yang muncul selalu membagi $n$.
Untuk tiap $d \mid n$, ada tepat $\varphi(d)$ pecahan berpenyebut $d$,
yaitu yang pembilangnya relatif prima dengan $d$.
Menghitung semua pecahan memberi $n = \sum_{d \mid n} \varphi(d)$.

\begin{example}
Verifikasikan identitas Gauss untuk $n = 12$.
\end{example}
\begin{proof}[Solusi]
Faktor dari $12$ adalah $1, 2, 3, 4, 6, 12$:
\[ \varphi(1) + \varphi(2) + \varphi(3) + \varphi(4) + \varphi(6) + \varphi(12)
= 1 + 1 + 2 + 2 + 2 + 4 = 12. \]
\end{proof}

\begin{fact}[$\varphi$ genap]
Untuk setiap $n \ge 3$, nilai $\varphi(n)$ genap.
\end{fact}

Pasangkan tiap $a$ yang relatif prima dengan $n$ bersama $n - a$,
yang juga relatif prima dengan $n$.
Pasangan ini selalu terdiri dari dua bilangan berbeda untuk $n \ge 3$:
jika $a = n - a$ maka $n = 2a$ dan $FPB(a, n) = a = 1$,
sehingga $n = 2$. Jadi $\varphi(n)$ genap.

\begin{example}
Buktikan tidak ada bilangan asli $n$ dengan $\varphi(n) = 7$.
\end{example}
\begin{proof}[Solusi]
Untuk $n \ge 3$, nilai $\varphi(n)$ selalu genap, jadi tidak mungkin $7$.
Tinggal periksa $\varphi(1) = 1$ dan $\varphi(2) = 1$.
\end{proof}

\section{Menjumlahkan atas Pembagi dan Fungsi M\"obius}

\begin{theorem}
Jika $f$ multiplikatif, maka $F(n) = \sum_{d \mid n} f(d)$ juga multiplikatif.
\end{theorem}

Ambil $m$ dan $n$ yang relatif prima.
Setiap faktor $d$ dari $mn$ tunggal berbentuk $d_1 d_2$ dengan $d_1 \mid m$
dan $d_2 \mid n$, dan $f(d_1 d_2) = f(d_1)f(d_2)$.
Maka jumlah atas semua $d$ memfaktorkan:
\[ F(mn) = \left(\sum_{d_1 \mid m} f(d_1)\right)\left(\sum_{d_2 \mid n} f(d_2)\right)
= F(m)F(n). \]

\begin{example}
Buktikan bahwa $\tau$ multiplikatif.
\end{example}
\begin{proof}[Solusi]
Tulis $\tau$ sebagai jumlah atas pembagi:
\[ \tau(n) = \sum_{d \mid n} 1. \]
Fungsi konstan $1$ jelas multiplikatif, jadi teorema di atas memberi
$\tau$ multiplikatif.
\end{proof}

\begin{definition}
Bilangan asli $n$ disebut \istilah{bebas kuadrat}{squarefree} apabila tidak ada
kuadrat prima yang membaginya. \istilah{Fungsi M\"obius}{M\"obius function} $\mu(n)$
didefinisikan sebagai
\[ \mu(n) = \begin{cases}
1 & \text{jika } n = 1, \\
(-1)^m & \text{jika $n$ squarefree dengan $m$ faktor prima}, \\
0 & \text{jika $n$ tidak squarefree}.
\end{cases} \]
\end{definition}

\begin{example}
Hitunglah $\mu(30)$.
\end{example}
\begin{proof}[Solusi]
Faktorisasi $30 = 2 \cdot 3 \cdot 5$: squarefree dengan $3$ faktor prima.
Jadi
\[ \mu(30) = (-1)^3 = -1. \]
\end{proof}

\begin{theorem}
Untuk setiap bilangan asli $n$ berlaku
\[ \sum_{d \mid n} \mu(d) = \begin{cases} 1 & \text{jika } n = 1, \\ 0 & \text{jika } n > 1. \end{cases} \]
\end{theorem}

Untuk $n = 1$ jelas. Untuk $n > 1$ dengan $m$ faktor prima berbeda,
hanya pembagi squarefree yang berkontribusi:
ada $\binom{m}{k}$ pembagi squarefree dengan $k$ faktor prima,
masing-masing bernilai $(-1)^k$. Jumlahnya:
\[ \sum_{k=0}^{m} \binom{m}{k}(-1)^k = (1 - 1)^m = 0. \]

\begin{example}
Verifikasikan teorema di atas untuk $n = 12$.
\end{example}
\begin{proof}[Solusi]
Faktor dari $12$ adalah $1, 2, 3, 4, 6, 12$:
\[ \mu(1) + \mu(2) + \mu(3) + \mu(4) + \mu(6) + \mu(12)
= 1 - 1 - 1 + 0 + 1 + 0 = 0. \]
\end{proof}

\begin{theorem}[Inversi M\"obius (\textit{M\"obius inversion})]
Misalkan $f$ dan $g$ fungsi aritmetika dengan $g(n) = \sum_{d \mid n} f(d)$
untuk setiap $n$. Maka
\[ f(n) = \sum_{d \mid n} \mu(d)\, g\left(\frac{n}{d}\right) \]
untuk setiap $n$.
\end{theorem}

\begin{proof}[Solusi]
Uraikan ruas kanan lalu tukar urutan penjumlahannya:
\begin{align*}
\sum_{d \mid n} \mu(d)\, g\left(\frac{n}{d}\right)
&= \sum_{d \mid n} \mu(d) \sum_{e \mid n/d} f(e) \\
&= \sum_{e \mid n} f(e) \sum_{d \mid n/e} \mu(d).
\end{align*}
Jumlah dalam bernilai $1$ bila $n/e = 1$ dan $0$ bila tidak,
jadi seluruhnya sama dengan $f(n)$.
\end{proof}

\begin{example}
Dengan inversi M\"obius pada identitas Gauss, hitunglah $\varphi(12)$.
\end{example}
\begin{proof}[Solusi]
Gauss mengatakan $n = \sum_{d \mid n} \varphi(d)$,
jadi inversi memberi $\varphi(n) = \sum_{d \mid n} \mu(d) \cdot \frac{n}{d}$.
Untuk $n = 12$:
\begin{align*}
\varphi(12) &= \mu(1) \cdot 12 + \mu(2) \cdot 6 + \mu(3) \cdot 4 \\
&\quad + \mu(4) \cdot 3 + \mu(6) \cdot 2 + \mu(12) \cdot 1 \\
&= 12 - 6 - 4 + 0 + 2 + 0 \\
&= 4.
\end{align*}
\end{proof}

\part*{Soal Latihan}

\section{Banyak Faktor}

\begin{enumerate}

	% Source: WMI 2019 Grade 5-6 (diterjemahkan)
	% Jawaban: 7
	\item (WMI 2019)
	Ada berapa bilangan asli antara $300$ sampai $600$ yang banyak faktor
	positifnya ganjil?

	% Source: AMO 2019 Grade 8 No. 18 (diterjemahkan)
	% Jawaban: 7
	\item (AMO 2019)
	Ada berapa bilangan $3$ digit yang tepat memiliki $3$ faktor positif?

	% Source: OSK SMA 2015 No. 1 (Isian Singkat)
	% Jawaban: 8
	\item (OSK 2015)
	Banyaknya faktor bulat positif dari 2015 adalah \dots

	% Source: OSK SMA 2008 No. 2 (Pilihan Ganda, disajikan sebagai isian)
	% Jawaban: 16
	\item (OSK 2008)
	Banyaknya faktor positif dari $5!$ adalah \dots

	% Source: OSK SMA 2007 No. 4 (Pilihan Ganda, disajikan sebagai isian)
	% Jawaban: 63
	\item (OSK 2007)
	Misalkan $H$ adalah himpunan semua faktor positif dari 2007.
	Banyaknya himpunan bagian dari $H$ yang tidak kosong adalah \dots

	% Source: OSK SMA 2012 No. 12 (Isian Singkat)
	% Jawaban: 420
	\item (OSK 2012)
	Ada berapa faktor positif dari $2^{7}3^{5}5^{3}7^{2}$
	yang merupakan kelipatan 6?

	% Source: OSK SMA 2011 No. 20 (Isian Singkat)
	% Jawaban: 378
	\item (OSK 2011)
	Ada berapa faktor positif dari $2^{7}3^{5}5^{3}7^{2}$
	yang merupakan kelipatan 10?

	% Source: OSK SMA 2010 No. 12 (Isian Singkat)
	% Jawaban: 16/81
	\item (OSK 2010)
	Misalkan $S$ menyatakan himpunan semua faktor positif dari $2010^{2}$.
	Sebuah bilangan diambil secara acak dari $S$.
	Peluang bilangan yang terambil habis dibagi 2010 adalah \dots

	% Source: OSK SMA 2017 Medan Ulang No. 5 (Isian Singkat)
	% Jawaban: 7^2016
	\item (OSK 2017)
	Suatu bilangan asli mempunyai 2017 pembagi positif yang salah satu
	di antaranya adalah 7. Bilangan yang dimaksud adalah \dots

	% Source: OSP SMA 2008 No. 1 (Isian Singkat)
	% Jawaban: 8
	\item (OSP 2008)
	Banyaknya pembagi positif dari 2008 adalah \dots

	% Source: AMC 10 2017 B (diterjemahkan, disajikan sebagai isian)
	% Jawaban: 1/19
	\item (AMC 10B 2017)
	Bilangan $21!$ memiliki lebih dari $60000$ faktor positif.
	Satu faktor dipilih secara acak.
	Peluang faktor yang terpilih ganjil adalah \dots

	% Source: AMC 10 2020 A (diterjemahkan, disajikan sebagai isian)
	% Jawaban: 23
	\item (AMC 10A 2020)
	Sebuah faktor positif dari $12!$ dipilih secara acak.
	Peluang faktor tersebut kuadrat sempurna (\textit{perfect square}) dapat ditulis $\frac{m}{n}$
	dalam bentuk paling sederhana. Nilai $m + n$ adalah \dots

	% Source: AMC 10 2019 A (diterjemahkan, disajikan sebagai isian)
	% Jawaban: 37
	\item (AMC 10A 2019)
	Ada berapa faktor positif dari $201^{9}$ yang merupakan kuadrat sempurna
	atau kubik sempurna (\textit{perfect cube}; atau keduanya)?

	% Source: AMC 12 2005 B (diterjemahkan, disajikan sebagai isian)
	% Jawaban: 2
	\item (AMC 12B 2005)
	Bilangan asli $n$ memiliki $60$ faktor positif dan $7n$ memiliki
	$80$ faktor positif. Bilangan bulat $k$ terbesar sehingga $7^{k}$
	membagi $n$ adalah \dots

	% Source: AIME 1988 No. 5 (diterjemahkan)
	% Jawaban: 634
	\item (AIME 1988)
	Peluang sebuah faktor positif dari $10^{99}$ yang dipilih secara acak
	merupakan kelipatan dari $10^{88}$ adalah $\frac{m}{n}$ dalam bentuk
	paling sederhana. Nilai $m + n$ adalah \dots

	% Source: AIME 1995 No. 6 (diterjemahkan)
	% Jawaban: 589
	\item (AIME 1995)
	Misalkan $n = 2^{31}3^{19}$. Ada berapa faktor positif dari $n^2$
	yang kurang dari $n$ tetapi tidak membagi $n$?

	% Source: OSN SMA 2004 No. 1 (Esai, disajikan sebagai isian)
	% Jawaban: 36 genap dan 12 ganjil
	\item (OSN 2004)
	Ada berapa pembagi genap dan pembagi ganjil dari $5^{6} - 1$?

	% Source: AIME 2005 I (diterjemahkan)
	% Jawaban: 25
	\item (AIME I 2005)
	Untuk bilangan asli $n$, notasi $\tau(n)$ menyatakan banyak faktor positif
	dari $n$. Definisikan $S(n) = \tau(1) + \tau(2) + \dots + \tau(n)$.
	Misalkan $a$ banyak bilangan asli $n \le 2005$ dengan $S(n)$ ganjil
	dan $b$ banyaknya dengan $S(n)$ genap. Nilai $|a - b|$ adalah \dots

	% Source: OSK SMA 2020 No. 6 (Isian Singkat)
	% Jawaban: 2
	\item (OSK 2020)
	Kuartet bilangan asli $(a, b, c, d)$ dikatakan keren jika memenuhi
	$b = a^{2} + 1$, $c = b^{2} + 1$, $d = c^{2} + 1$ dan
	$\tau(a) + \tau(b) + \tau(c) + \tau(d)$ bilangan ganjil.
	Banyaknya kuartet keren $(a, b, c, d)$ dengan $a, b, c, d < 10^{6}$
	adalah \dots\ (Catatan: $\tau(k)$ menyatakan banyaknya faktor positif dari $k$.)

	% Source: AIME 2019 I (diterjemahkan)
	% Jawaban: 540
	\item (AIME I 2019)
	Notasi $\tau(n)$ menyatakan banyak faktor positif dari $n$.
	Tentukan jumlah dari enam bilangan asli terkecil $n$ yang memenuhi
	$\tau(n) + \tau(n + 1) = 7$.

\end{enumerate}

\section{Jumlah Faktor}

\begin{enumerate}[resume]

	% Source: AMC 10 2008 A (diterjemahkan, disajikan sebagai isian)
	% Jawaban: 6
	\item (AMC 10A 2008)
	Untuk bilangan asli $n$, notasi $\langle n \rangle$ menyatakan jumlah
	semua faktor positif dari $n$ kecuali $n$ sendiri.
	Sebagai contoh, $\langle 4 \rangle = 1 + 2 = 3$ dan
	$\langle 12 \rangle = 1 + 2 + 3 + 4 + 6 = 16$.
	Nilai dari $\langle \langle \langle 6 \rangle \rangle \rangle$ adalah \dots

	% Source: AMC 10 2024 B (diterjemahkan, disajikan sebagai isian)
	% Jawaban: 6
	\item (AMC 10B 2024)
	Misalkan $N$ adalah hasil kali semua faktor positif dari $42$.
	Digit satuan dari $N$ adalah \dots

	% Source: OSK SMA 2023 No. 16 (Isian Singkat)
	% Jawaban: 27
	\item (OSK 2023)
	Misalkan $N = 2^{a} \cdot 3^{b}$ dengan $a, b$ bilangan asli.
	Jika hasil kali semua faktor positif dari $N$ adalah $24^{60}$,
	maka nilai $ab$ adalah \dots

	% Source: AMC 12 2008 B (diterjemahkan, disajikan sebagai isian)
	% Jawaban: 11
	\item (AMC 12B 2008)
	Jumlah logaritma basis $10$ dari semua faktor positif dari $10^{n}$
	adalah $792$. Nilai $n$ adalah \dots

	% Source: AIME 1986 No. 8 (diterjemahkan)
	% Jawaban: 141
	\item (AIME 1986)
	Misalkan $S$ adalah jumlah logaritma basis $10$ dari semua faktor
	positif sejati dari $1000000$ (yaitu semua faktor kecuali bilangan itu
	sendiri). Bilangan bulat terdekat ke $S$ adalah \dots

	% Source: AMC 10 2021 B (diterjemahkan, disajikan sebagai isian)
	% Jawaban: 1/14
	\item (AMC 10B 2021)
	Misalkan $N = 34 \cdot 34 \cdot 63 \cdot 270$.
	Perbandingan jumlah faktor ganjil dari $N$ terhadap jumlah faktor genap
	dari $N$ adalah \dots

	% Source: AMC 12 2021 Fall B (diterjemahkan, disajikan sebagai isian)
	% Jawaban: 1/192
	\item (AMC 12 2021 Fall B)
	Untuk bilangan asli $n$, misalkan $f(n)$ adalah hasil bagi jumlah semua
	faktor positif dari $n$ oleh $n$.
	Sebagai contoh, $f(14) = (1 + 2 + 7 + 14) \div 14 = \frac{12}{7}$.
	Nilai dari $f(768) - f(384)$ adalah \dots

	% Source: AIME 2000 I (diterjemahkan)
	% Jawaban: 248
	\item (AIME I 2000)
	Misalkan $S$ adalah jumlah semua bilangan berbentuk $\frac{a}{b}$ dengan
	$a$ dan $b$ faktor positif dari $1000$ yang relatif prima (\textit{coprime}).
	Bilangan bulat terbesar yang tidak melebihi $\frac{S}{10}$ adalah \dots

	% Source: OSK SMA 2026 No. 15 (Isian Singkat)
	% Jawaban: 2010
	\item (OSK 2026)
	Misalkan $\sigma(n)$ adalah jumlah faktor positif dari $n$.
	Jumlah semua bilangan asli $n \le 225$ sehingga $\sigma(n)$ adalah
	bilangan ganjil adalah \dots

	% Source: AIME 2021 I (diterjemahkan)
	% Jawaban: 125
	\item (AIME I 2021)
	Notasi $\sigma(a)$ menyatakan jumlah faktor positif dari $a$.
	Misalkan $n$ bilangan asli terkecil sehingga $\sigma(a^{n}) - 1$
	habis dibagi $2021$ untuk setiap bilangan asli $a$.
	Jumlah faktor prima dalam faktorisasi prima dari $n$ adalah \dots

	% Source: Modul Pra-OSP FungsiAritmetika (juga di Khurmi Bab 3 Latihan 3.7.5)
	\item (Pra-OSP Fungsi Aritmetika)
	Buktikan bahwa untuk setiap bilangan asli $n$ berlaku
	\[ \frac{\sigma(n)}{\tau(n)} \ge \sqrt{n}. \]

\end{enumerate}

\section{Fungsi Totient}

\begin{enumerate}[resume]

	% Source: AMO 2022 Grade 9 No. 20 (diterjemahkan)
	% Jawaban: 24
	\item (AMO 2022)
	Untuk bilangan asli $n$, notasi $\phi(n)$ menyatakan banyak bilangan asli
	yang kurang dari $n$ dan relatif prima dengan $n$.
	Nilai dari $\phi(35)$ adalah \dots

	% Source: OSP SMA 2008 No. 14 (Isian Singkat)
	% Jawaban: 10
	\item (OSP 2008)
	Banyak bilangan bulat positif $n$ kurang dari 2008 yang mempunyai tepat
	$\frac{n}{2}$ bilangan kurang dari $n$ dan relatif prima terhadap $n$
	adalah \dots

	% Source: Komposisi sendiri (latihan klasik)
	% Jawaban: 7, 9, 14, 18
	\item
	Tentukan semua bilangan asli $n$ sehingga $\varphi(n) = 6$.

	% Source: Modul Pra-OSP FungsiAritmetika
	\item (Pra-OSP Fungsi Aritmetika)
	Buktikan bahwa ada tak hingga bilangan asli $n$ yang memenuhi
	$\varphi(n) = \frac{n}{3}$.

\end{enumerate}

\section{Campuran}

\begin{enumerate}[resume]

	% Source: OSK SMA 2024 No. 13 (Isian Singkat)
	% Jawaban: 90
	\item (OSK 2024)
	Untuk setiap bilangan asli $n$, misalkan $f(n)$ menyatakan faktor ganjil
	terbesar dari $n$ dan $p(n) = f(n) + f(n+1) + \dots + f(2n)$.
	Jika $p(n) = 8145$, maka nilai dari $n$ adalah \dots

	% Source: AIME 1987 No. 3 (diterjemahkan)
	% Jawaban: 182
	\item (AIME 1987)
	Faktor sejati (\textit{proper divisor}) dari bilangan asli adalah faktor positif selain $1$ dan
	bilangan itu sendiri. Bilangan asli lebih dari $1$ disebut \textit{nice}
	jika sama dengan hasil kali semua faktor sejatinya yang berbeda.
	Jumlah sepuluh bilangan \textit{nice} pertama adalah \dots

	% Source: AIME 2005 I (diterjemahkan)
	% Jawaban: 109
	\item (AIME I 2005)
	Ada berapa bilangan asli yang tepat memiliki tiga faktor sejati
	(faktor positif kecuali bilangan itu sendiri),
	yang masing-masing kurang dari $50$?

	% Source: AMC 10 2016 A (diterjemahkan, disajikan sebagai isian)
	% Jawaban: 325
	\item (AMC 10A 2016)
	Untuk suatu bilangan asli $n$, bilangan $110n^{3}$ memiliki $110$ faktor
	positif. Ada berapa faktor positif dari bilangan $81n^{4}$?

	% Source: AIME 2025 II (diterjemahkan)
	% Jawaban: 237
	\item (AIME II 2025)
	Misalkan $A$ himpunan semua faktor positif dari $2025$.
	Dipilih secara acak satu himpunan bagian $B$ dari $A$.
	Peluang $B$ merupakan himpunan tak kosong yang $KPK$ semua anggotanya
	sama dengan $2025$ adalah $\frac{m}{n}$ dalam bentuk paling sederhana.
	Nilai $m + n$ adalah \dots

	% Source: OSK SMA 2024 No. 9 (Isian Singkat)
	% Jawaban: 92
	\item (OSK 2024)
	Misalkan $k$ adalah bilangan bulat positif terkecil kelipatan 2024
	yang memiliki 28 faktor positif.
	Sisa hasil bagi $k$ oleh 100 adalah \dots

	% Source: AIME 1990 No. 5 (diterjemahkan)
	% Jawaban: 432
	\item (AIME 1990)
	Misalkan $n$ bilangan asli terkecil yang merupakan kelipatan $75$
	dan memiliki tepat $75$ faktor positif.
	Nilai dari $\frac{n}{75}$ adalah \dots

	% Source: OSP SMA 2009 No. 18 (Isian Singkat)
	% Jawaban: 2
	\item (OSP 2009)
	Misalkan $n$ bilangan asli terkecil yang mempunyai tepat 2009 faktor
	dan $n$ merupakan kelipatan 2009.
	Faktor prima terkecil dari $n$ adalah \dots

	% Source: AIME 2026 I (diterjemahkan)
	% Jawaban: 244
	\item (AIME I 2026)
	Misalkan $N$ adalah banyak faktor positif dari $17017^{17}$ yang bersisa
	$5$ jika dibagi $12$. Sisa $N$ jika dibagi $1000$ adalah \dots

\end{enumerate}

\section{Tantangan}

\begin{enumerate}[resume]

	% Source: Kanada MO 1999 (via Modul Pra-OSP FungsiAritmetika)
	% Jawaban: 1 dan 9
	\item (Kanada MO 1999)
	Tentukan semua bilangan asli $n$ yang memenuhi $n = \tau(n)^2$.

	% Source: Modul Pra-OSP FungsiAritmetika
	% Jawaban: 1, 3, 8, 10, 18, 24, 30
	\item (Pra-OSP Fungsi Aritmetika)
	Carilah semua solusi persamaan $\varphi(n) = \tau(n)$.

	% Source: OSN SMA 2010 No. 6 (Esai)
	% Jawaban: 4 atau bilangan prima
	\item (OSN 2010)
	Cari semua bilangan asli $n > 1$ sedemikian sehingga
	\[ \tau(n) + \varphi(n) = n + 1. \]
	Dalam hal ini, $\tau(n)$ menyatakan banyaknya bilangan asli yang habis
	membagi $n$, dan $\varphi(n)$ menyatakan banyaknya bilangan asli yang
	kurang dari $n$ dan relatif prima terhadap $n$.

	% Source: Filipina MO 2017 (via Modul Pra-OSP FungsiAritmetika)
	% Jawaban: kesamaan tercapai jika dan hanya jika n = 1
	\item (Filipina MO 2017)
	Buktikan bahwa
	\[ \frac{1}{\sigma(n)} + \frac{1}{\varphi(n)} \ge \frac{2}{n} \]
	dan tentukan kapan kesamaan tercapai.

	% Source: Turki MO 2022 (via Modul Pra-OSP FungsiAritmetika)
	\item (Turki MO 2022)
	Buktikan bahwa untuk setiap bilangan asli $k, n$ yang memenuhi
	$k \ge n!$ berlaku $\varphi(k) \ge (n - 1)!$.

	% Source: Irlandia MO 1998 (via Khurmi Bab 3 Latihan 3.7.10)
	% Jawaban: 1998 dan 3834
	\item (Irlandia MO 1998)
	Tentukan semua bilangan asli $d$ yang memiliki tepat $16$ faktor positif
	$d_1 < d_2 < \dots < d_{16}$ sehingga $d_6 = 18$ dan $d_9 - d_8 = 17$.

\end{enumerate}

\section{Referensi}
\begin{itemize}
    \item Aditya Khurmi, \textit{Modern Olympiad Number Theory},
    Bab 3 (Fungsi Aritmetika), 2020.
    \item Modul Pra-OSP: Fungsi Aritmetika (Azzam dan Akeyla).
    \item Unit Modulo dan Keterbagian (rumus $\varphi$ dan teorema Euler).
\end{itemize}

\end{document}
